% ==============================================================================
% Section 2: Per-Layer Spatial Probes (the diagnostic lens; also the readout head)
% ==============================================================================

\section{Per-Layer Spatial Probes}
\label{sec:probe}
\label{sec:method:stage1}

We diagnose \emph{where} GUI grounding signals reside across transformer depth by attaching one lightweight probe per layer to a frozen backbone. The same probe set later forms the readout head of the retrofit (\S\ref{sec:method}). Probes introduce no new spatial knowledge: the backbone is frozen, and the gradient from layer $\ell$ reaches only $\phi_\ell$.

\paragraph{Prefill-forced inference.}
Coordinate generation serializes a location into up to eight number tokens, but the spatial intent is already encoded before decoding begins (\S\ref{sec:analysis}). We prefill a fixed template ending in an anchor token, $\texttt{click(}\,\texttt{<|ground|>}\,\texttt{<|pointer\_start|>}\cdots\texttt{)}$, and run a single forward pass. At the anchor position the model has processed the full screenshot and instruction but no coordinate value, so there is zero label leakage.

\paragraph{Cross-attention probe.}
Let $f_\theta$ be a frozen $L$-layer backbone. For each probed layer $\ell\in\mathcal{L}_{\mathrm{probe}}$, drawn from the deeper half of the network where spatial signal concentrates, we read the anchor query $\mathbf{h}^{(\ell)}_{t^*}\in\mathbb{R}^d$ at the \texttt{<|ground|>} position and visual keys $\mathbf{H}^{(\ell)}_v\in\mathbb{R}^{N_v\times d}$ over the $N_v$ patch tokens. Each probe $\phi_\ell$ computes a patch-level posterior via multi-head cross-attention with a learnable head gate $\boldsymbol{\alpha}$:
\begin{equation}
\begin{aligned}
  \mathbf{Q}_\ell &= W_q^{(\ell)}\,\mathrm{Norm}(\mathbf{h}^{(\ell)}_{t^*}),\quad
  \mathbf{K}_\ell = W_k^{(\ell)}\,\mathrm{Norm}(\mathbf{H}^{(\ell)}_v),\\[-2pt]
  p_\ell &= \sigma\!\left(\textstyle\sum_h \sigma(\boldsymbol{\alpha})_h\,\frac{\mathbf{Q}_{\ell,h}^{\!\top}\mathbf{K}_{\ell,h}}{\sqrt{d_h}}\right)\in\Delta^{N_v-1},
\end{aligned}
\label{eq:probe}
\end{equation}
where $\mathrm{Norm}(\cdot)$ is the RMS pre-scale followed by LayerNorm, and $W_q^{(\ell)},W_k^{(\ell)}$ are full-rank projections so each probe learns a geometry suited to its layer. Each probe is ${\approx}4.2$M parameters ($d{=}4096$, $n_h{=}8$, $d_h{=}64$); the full set totals ${\approx}42$M, or $0.52\%$ of a 7B backbone.

\paragraph{RMSNorm pre-scale.}
Hidden-state $\ell_2$ norms grow with depth, geometrically misaligning layers. We pre-scale $\mathbf{h}\leftarrow\mathbf{h}/(\|\mathbf{h}\|_2/\sqrt{d})$ before each LayerNorm, bounding the input norm to ${\approx}\sqrt{d}$ regardless of depth and preventing \texttt{bfloat16} overflow.

\paragraph{Supervision.}
Each probe is supervised with an anisotropic Gaussian centered at the bbox centroid:
\begin{equation}
  y_i \propto \exp\!\left(-\tfrac{(x_i-c_x)^2}{2\sigma_x^2} - \tfrac{(y_i-c_y)^2}{2\sigma_y^2}\right),\quad \sigma_x=\eta w_b,\ \sigma_y=\eta h_b
  \label{eq:gaussian}
\end{equation}
which reduces over-penalization from patch-grid quantization on small targets (\S\ref{sec:analysis}). Stage~1 minimizes $\mathcal{L}_{\mathrm{S1}}=\tfrac{1}{|\mathcal{L}_{\mathrm{probe}}|}\sum_\ell D_{\mathrm{KL}}(\mathbf{y}\,\|\,p_\ell)$ with the backbone receiving zero gradient. Training-data construction is detailed in the supplementary material.
